\section{Installing the runtime}
\label{sec:runtime}

\subsection{Staging the binary}
Because the container has no downloader and no root, the runtime is fetched on
the developer VM and staged through \cmd{/shared}.

\begin{lstlisting}
# on the developer VM
mkdir -p /shared/$USER/opt /shared/$USER/ollama-models
cd /shared/$USER/opt
curl -sL -o ollama.tar.zst \
  https://github.com/ollama/ollama/releases/download/v0.34.0/ollama-linux-amd64.tar.zst
mkdir -p ollama && zstd -dc ollama.tar.zst | tar -xf - -C ollama
\end{lstlisting}

This yields \cmd{bin/ollama} plus the bundled CUDA runners in \cmd{lib/ollama}.
No installation step, no systemd unit, no root required --- which is precisely
why it works in a container where none of those are available.

\subsection{Where Ollama sits}
Ollama is a management layer, not the inference engine and not the model.

\begin{figure}[!ht]
  \centering
  \begin{tikzpicture}[
    font=\small,
    lay/.style={draw=atBorderStrong, minimum width=10.5cm, minimum height=0.95cm,
                align=center},
    arr/.style={{Latex[length=2mm]}-{Latex[length=2mm]}, draw=atFgMuted},
  ]
  \node[lay, fill=atBg]       (cli) {client: \cmd{curl}, editor plugin, OpenAI SDK};
  \node[lay, fill=atBgSubtle, below=2mm of cli] (oll)
    {\textbf{Ollama} --- process supervision, model registry,\\
     prompt templating, HTTP API};
  \node[lay, fill=atBgMuted, below=2mm of oll] (eng)
    {\cmd{llama-server} (llama.cpp / ggml) + CUDA backend};
  \node[lay, fill=atBgSubtle, below=2mm of eng] (wts)
    {weights: GGUF on \cmd{/shared} (Table~\ref{tab:gpu-model-inventory})};
  \draw[arr] (cli) -- (oll);
  \draw[arr] (oll) -- (eng);
  \draw[arr] (eng) -- (wts);
  \end{tikzpicture}
  \caption{The serving stack, from client call down to the weights on disk.}
  \label{fig:gpu-ollama-stack}
\end{figure}

Ollama supervises a separate \cmd{llama-server} process; llama.cpp and its CUDA
kernels perform the arithmetic. Ollama contributes four things: it fetches and
stores weights as blobs and manifests, it loads and unloads models with a
keep-alive timer, it renders chat turns into the exact token format the model
was trained on, and it exposes both a native API and an OpenAI-compatible one.
It also sizes the context window to measured VRAM at start-up.

For a single user on a single card this is the right trade. For a shared,
heavily concurrent endpoint, vLLM's continuous batching would serve better at
the cost of a much heavier install.
